%%%
%
% $Autor: Siddhanth Sharma $
% $Datum: 2026-01-02 $
%
% !TeX spellcheck = GB
% !TeX program = pdflatex
% !TeX encoding = utf8
%
%%%

\chapter{Development Environment -- Python}
\index{Domain Knowledge -- Tools}

\noindent
This chapter documents the \textbf{project-relevant} tools and the system environment used to build and run the license plate detection system.
The emphasis is on \textbf{usage, configuration, and validation} required for reproducible execution (not on tutorials).

\section{Version}
\noindent
Python is the primary runtime used to execute the end-to-end pipeline (camera capture, image preprocessing, model inference, OCR, and logging).
The interpreter version was verified directly on the target device using \SHELL{python3 --version}.

\vspace{1mm}
\noindent
\textbf{Recorded value:} Python 3.6.9 (\FILE{/usr/bin/python3}).

\section{Description}
\noindent
Python is used for the complete processing chain:
\begin{itemize}
	\item \textbf{Camera access and image acquisition} (USB camera frames via OpenCV).
	\item \textbf{Image preprocessing} (cropping, resizing, binarization and contrast adjustment for OCR).
	\item \textbf{Detection inference} (license plate detection model executed on the Jetson).
	\item \textbf{OCR integration} (text extraction from the cropped plate region).
	\item \textbf{Local logging} (timestamps and detection results persisted locally).
\end{itemize}

\vspace{1mm}
\noindent
The following Python modules are relevant for this project:
\begin{itemize}
	\item \textbf{OpenCV (cv2):} camera access, preprocessing, and visualization.
	\item \textbf{NumPy:} array operations and numerical processing.
	\item \textbf{SQLite (sqlite3):} lightweight, file-based storage for entries, exits, and timestamps.
	
	\item \textbf{Detection runtime:}
	\begin{itemize}
		\item \textbf{Optimized deployment:}  
		TensorRT engine loaded by the Flask dashboard.
		
		{\sloppy\small
			\begin{itemize}
				\item Engine file: \FILE{/media/myssd/jetson-inference/data/networks/best.engine}
				\item Loader script: \FILE{Code/LicencePlateDetection/jetson-inference/python/www/flask/dashboard2.py}
			\end{itemize}
		}
		
		\item \textbf{Prototype / runtime fallback:}  
		Ultralytics YOLO (PyTorch) loading trained weights.
		
		{\sloppy\small
			\begin{itemize}
				\item Weights file: \FILE{best.pt}
				\item Detector module: \FILE{Code/LicencePlateDetection/jetson\_app/detector.py}
				\item Prototype dashboard: \FILE{Code/LicencePlateDetection/jetson\_app/dashboard.py}
			\end{itemize}
		}
	\end{itemize}
	
	\item \textbf{OCR runtime:}  
	EasyOCR (\texttt{easyocr==1.7.1}).  
	In the deployment dashboard (\FILE{dashboard2.py}), OCR is executed on the CPU
	(\PYTHON{gpu=False}) to reduce GPU memory pressure.
	
	\item \textbf{UI layer (final demo):}
	\begin{itemize}
		\item \textbf{Web dashboard (Jetson):}  
		Flask application started via \SHELL{python3 dashboard2.py}, listening on
		\texttt{0.0.0.0:8055}.
		
		{\sloppy\small
			\begin{itemize}
				\item Application path: \FILE{Code/LicencePlateDetection/jetson-inference/python/www/flask/dashboard2.py}
				\item Endpoints:
				\begin{itemize}
					\item \texttt{/} — Web UI
					\item \texttt{/video\_feed} — MJPEG video stream
					\item \texttt{/api/billing} — JSON billing API
				\end{itemize}
			\end{itemize}
		}
		
		\item \textbf{Prototype UI (not used in final demo):}  
		Flet-based kiosk dashboard implemented as a GUI application.
		
		{\sloppy\small
			\begin{itemize}
				\item Path: \FILE{Code/LicencePlateDetection/jetson\_app/dashboard.py}
			\end{itemize}
		}
	\end{itemize}
\end{itemize}

\section{Installation}
\noindent
Python is provided by the Jetson OS image (JetPack/L4T) and does not require separate installation.
Project dependencies were installed via \SHELL{pip}/\SHELL{pip3}, ideally from a pinned requirements file.

\vspace{1mm}
\noindent
\textbf{Reproducibility items to document:}
\begin{itemize}
	\item \textbf{Requirements file:} \FILE{Code/LicencePlateDetection/requirements.txt}.
	\item \textbf{Installation command:} from the repository root, \SHELL{python3 -m pip install -r Code/LicencePlateDetection/requirements.txt}.
	\item \textbf{Virtual environment:} the development repository contains a venv folder at \FILE{Code/LicencePlateDetection/.venv}. If used on the Jetson, it is activated with \SHELL{source Code/LicencePlateDetection/.venv/bin/activate} (created via \SHELL{python3 -m venv .venv}).
\end{itemize}

\vspace{1mm}
\noindent
\textbf{Recorded runtime packages (Jetson demo device):}
\begin{itemize}
	\item \textbf{PyTorch:} \texttt{1.10.0a0+git36449ea} (CUDA available: \texttt{True}).
	\item \textbf{TensorRT (Python):} \texttt{tensorrt==8.2.1.8}.
	\item \textbf{EasyOCR:} \texttt{easyocr==1.7.1}.
	\item \textbf{Flask:} \texttt{Flask==2.0.3}.
\end{itemize}

\section{Configuration}
\noindent
Only project-relevant configuration is required:
\begin{itemize}
		\item \textbf{Interpreter selection:} system Python (\FILE{/usr/bin/python3}) vs.\ virtual environment activation (\SHELL{source Code/LicencePlateDetection/.venv/bin/activate}).
		\item \textbf{Hardware acceleration visibility:} CUDA is available from Python (\PYTHON{torch.cuda.is\_available()} $\rightarrow$ \texttt{True}).
		\item \textbf{Memory-aware runtime setup:} OCR is configured to run on CPU (\PYTHON{gpu=False}) to keep GPU memory available for TensorRT inference and reduce OOM risk on the 4\,GB unified-memory device.
		\item \textbf{Post-processing:} Non-Maximum Suppression (NMS) is performed inside the YOLO runtime. In our scripts we control this via the model thresholds (e.g., \PYTHON{conf=0.2} and \PYTHON{iou=0.5} in \FILE{Code/LicencePlateDetection/jetson\_app/detector.py}).
		\item \textbf{Offline OCR models:} \FILE{Code/LicencePlateDetection/jetson\_app/ocr\_pipeline.py} disables automatic model downloads (\PYTHON{download\_enabled=False}); therefore the EasyOCR model files must already be present on the Jetson (downloaded once during initial setup).
\end{itemize}

\section{First Steps}
\noindent
Before running the full pipeline, the environment was validated with a short checklist:
\begin{itemize}
	\item \textbf{Interpreter check:} Python starts and reports the expected version.
	\item \textbf{Imports:} \PYTHON{cv2}, \PYTHON{numpy}, \PYTHON{tensorrt}, \PYTHON{pycuda}, \PYTHON{easyocr}, \PYTHON{flask} import without errors.
	\item \textbf{Camera access:} a camera device can be opened from Python and returns a frame.
	\item \textbf{Model availability:} the final model/engine file exists at the expected path and can be loaded.
\end{itemize}

\noindent
Figure~\ref{fig:python_verification_flow} summarizes the Python verification workflow used before executing the main pipeline.

\begin{figure}[h]
	\centering
	\begin{tikzpicture}[
		node distance=10mm,
		box/.style={draw, rounded corners, align=center, inner sep=6pt, text width=0.78\linewidth},
		arrow/.style={-Latex, thick}
		]
		\node[box] (jetson) {Jetson Nano system (JetPack/L4T installed)};
		\node[box, below=of jetson] (python) {Python Interpreter (\texttt{python3})};
		\node[box, below=of python] (deps) {Dependencies installed (\texttt{pip}/\texttt{pip3})\\(OpenCV, NumPy, detection/OCR runtime, SQLite)};
		\node[box, below=of deps] (hello) {Verification script\\(\texttt{python\_project\_sanity\_check.py})};
		\node[box, below=of hello] (ok) {Success\\(imports OK, camera frame read OK)};

		\draw[arrow] (jetson) -- (python);
		\draw[arrow] (python) -- (deps);
		\draw[arrow] (deps) -- (hello);
		\draw[arrow] (hello) -- (ok);
	\end{tikzpicture}
	\caption{Python environment verification flow on the Jetson device.}
	\label{fig:python_verification_flow}
\end{figure}

\section{Program ``Hello World''}
\noindent
This is not a generic \PYTHON{print("Hello World")} example.
Instead, a minimal project-relevant script validates that the Python environment can (i) import the required libraries and (ii) access the camera to read a frame.
Optionally, it can be extended to load the final detection runtime (model/engine) once the final deployment choice is fixed.

\lstinputlisting[
language=Python,
caption={Project-relevant Python ``Hello World'' (imports + camera frame read)},
label={lst:python_hello_world_project},
basicstyle=\ttfamily\small,
keywordstyle=\color{blue},
commentstyle=\color{gray},
stringstyle=\color{green!40!black},
breaklines=true,
showstringspaces=false,
numbers=left,
numberstyle=\tiny\color{gray},
frame=single,
rulecolor=\color{black!30}
]{../../Code/SystemTools/python_project_sanity_check.py}


